\section{Simulation Study: Controlled Evidence Degradation (RQ3, RQ4)}
\label{sec:simulation}

We now exercise the framework in the one setting where ground truth is available
by construction: a synthetic estate whose complete history is generated, then
degraded under experimental control. The study is executed, not merely designed;
every number below was produced by the released artifact ($\approx$560 lines of
dependency-free Python), is deterministic given the published seeds, and the
results table is generated programmatically from the run outputs. Its epistemic
role is stated with care, because a simulation whose apparatus is written by the
authors of the model cannot ``confirm'' claims that the apparatus embodies.
We therefore separate, explicitly, (i)~\emph{demonstrations}: behaviors entailed
by design choices that mirror real toolchains, whose value is quantification and
whose evidential force is tested by \emph{removal arms} in which the choice is
switched off; and (ii)~\emph{falsifiable tests}: predictions that the apparatus
does not build in and that the data could have refuted. The study validates
construct computability end to end---including $\ED(t)$ itself
(\cref{sec:sim-workload})---and estimates no real-world magnitudes.

\subsection{Apparatus}
\label{sec:sim-apparatus}

\begin{figure}[t]
\centering
\resizebox{\columnwidth}{!}{% Degradation experiment pipeline (single column, resized)
\begin{tikzpicture}[node distance=3.6mm and 5mm]
  \node[stage, minimum width=24mm] (gen)
    {\textbf{generator}\\{\smlbl 8 services $\cdot$ 6 stores}\\[-1pt]{\smlbl 6 links/chain $\cdot$ redundancy}\\[-1pt]{\smlbl sparse / dense profiles}};
  \node[brok, right=of gen, minimum width=24mm] (deg)
    {\textbf{degrader}\\{\smlbl intent $\cdot$ approval $\cdot$ corr-ids}\\[-1pt]{\smlbl records $\cdot$ timestamps $\cdot$ ids+time}\\[-1pt]{\smlbl all-classes; $p \in \{1,5,10,20,40\}\%$}};
  \node[stage, right=of deg, minimum width=24mm] (rec)
    {\textbf{reconstructor}\\{\smlbl exact joins}\\[-1pt]{\smlbl redundant paths}\\[-1pt]{\smlbl heuristic linkage}};
  \draw[flow] (gen) -- node[lbl, above]{$\Hist,\Corp$} (deg);
  \draw[flow] (deg) -- node[lbl, above]{$\Corp'$} (rec);

  \node[attn, below=4.5mm of gen, minimum width=24mm] (truth)
    {ground truth $\Hist$\\{\smlbl retained for scoring}};
  \node[evid, below=4.5mm of rec, minimum width=24mm] (score)
    {\textbf{scoring}\\{\smlbl coverage $\cdot$ false joins}\\[-1pt]{\smlbl irrecoverable $\cdot$ effort}};
  \draw[eflow] (gen) -- (truth);
  \draw[flow] (rec) -- (score);
  \draw[eflow] (truth) -- (score);
  \node[lbl, below=1.5mm of deg, xshift=2mm, text=black!60]
    {20 seeds $\times$ 2 profiles $\times$ 7 arms $\times$ 6 levels};
\end{tikzpicture}
}
\caption{The evidence-degradation experiment. A generator emits a complete
synthetic estate (ground-truth chains and their records across six stores); a
degrader removes a controlled fraction $p$ of one or more evidence classes; an
algorithmic reconstructor rebuilds each chain from surviving records (exact key
joins, then redundant-key paths, then heuristic linkage); scoring against ground
truth yields coverage, false joins, two irrecoverability constructs, invariant
effort, and $\ED$ under a declared workload.}
\label{fig:experiment}
\end{figure}

\subsubsection{Ground truth}
Each run generates an estate of eight services whose changes produce six-record
chains across six stores---ticket (ITSM), commit (forge), approval (review),
build (CI), artifact (registry), deployment (CD), plus the runtime configuration
record serving as query entry point---linked by the keys real toolchains use.
Six ground-truth links per chain instantiate the six interrogatives. The
generator includes \emph{redundancy} (the artifact digest recorded in both build
and configuration records), mirroring practice; because this choice matters, a
removal arm disables it (\cref{sec:sim-sensitivity}). Inter-record time gaps are
drawn below the reconstructor's default heuristic window; a removal arm breaks
that relationship too.

\subsubsection{Factors and arms}
The main factorial crosses \emph{degradation class}---intent links, approval
links, correlation identifiers, whole-record deletion, timestamps, the
\emph{pairwise} arm (correlation identifiers $+$ timestamps, which instantiates
the premise of \cref{prop:superadd}: neither class alone destroys records), and
an \emph{all-classes} arm---with \emph{level}
$p \in \{0,1,5,10,20,40\}\%$ and \emph{activity profile} (sparse: 50
chains/service, 180 days, 12 actors; dense: 150 chains/service, 90 days, 6
actors), 20 seeds per cell. Sensitivity arms (\cref{sec:sim-sensitivity}) vary
one design choice at a time: heuristic window below the maximum true gap;
redundancy off; ticket filer $\ne$ commit author for 30\% of chains; and a
single-knob density sweep (chains/service $\in \{50,100,150\}$ with actor pool
and window fixed), the latter because the sparse/dense profiles deliberately
bundle three knobs and cannot isolate density.

\subsubsection{Reconstructor and invariant effort accounting}
For each configuration record, the reconstructor rebuilds the chain backward:
surviving exact keys and redundant key paths first (implemented as index
lookups and charged \emph{zero} effort, so that effort does not depend on
incidental data-structure choices), then heuristic linkage (same-service
candidates within a 48-hour window, actor agreement where applicable; unique
candidate accepted; multiple candidates resolved nearest-in-time---a realistic
policy that can produce \emph{false joins}). Effort is counted only at
heuristic decision points, as two implementation-invariant quantities: the
number of links requiring non-exact resolution (\emph{fallback links}) and the
candidate records examined there.

\subsubsection{Scoring: two irrecoverability constructs}
Against ground truth we score \emph{link coverage}, \emph{false-join rate}
(incorrect/accepted), \emph{chain coverage}, and two distinct irrecoverability
constructs whose separation \cref{sec:model-graph} demands: \emph{record loss}
(a required record exists in no store---the destroyed-source clause of
\cref{def:irrec}) and \emph{identifiability-based irrecoverability}
(\cref{def:irrec} in full: record destroyed, \emph{or} no key path survives and
the discriminators---timestamps, actor, service, window---fail to single out
the true parent). The second is decidable here because the closed world knows
the truth; it is the construct the definition means, and it behaves very
differently from record loss (\cref{fnd:constructs}). One scoring asymmetry is
noted for reusers: five links are scored on the recovered parent's key, while
the approval link---whose evidential record sits on the child side---is scored
on the approval record's key; the artifact documents the orientation.

\subsubsection{$\ED$ under a declared workload}
\label{sec:sim-workload}
To instantiate \cref{def:debt} end to end, we declare a synthetic workload: one
full-chain audit query per running configuration, unit weight; resolution cost
$= 1$ unit per fallback link $+\,0.01$ per candidate examined; write-off
penalty $\pi = 50$ units per identifiability-irrecoverable link. The complete
corpus $\Corp^{*}$ resolves every link exactly, so its cost is zero and
measured cost is excess cost. The declarations are arbitrary in magnitude and
declared precisely because \cref{def:debt} requires them---the point is the
demonstrated pipeline from census to a computed $\ED(t)$, and the sensitivity
of $\ED$'s composition to the declarations is itself a result
(\cref{fnd:ed}).

\subsection{Demonstrations, Predictions, and Results}
\label{sec:sim-results}

\begin{table*}[t]
\centering
\caption{Selected Results, Generated Programmatically from the Run Outputs\\(Mean over 20 Seeds; SD in Parentheses). Cov = Link Coverage,\\FJ = False-Join Rate, Irr = Identifiability-Based Irrecoverable Chains,\\FB = Fallback Links per Chain, ED = Computed Debt (Units, \Cref{sec:sim-workload})}
\label{tab:results}
\footnotesize
\begin{tabular}{@{}llrlllll@{}}
\toprule
\textbf{Class} & \textbf{Prof.} & $p$ & \textbf{Cov} & \textbf{FJ} & \textbf{Irr} & \textbf{FB} & \textbf{ED} \\
\midrule
combined & sparse & 0\% & 1.000 (0.000) & 0.000 (0.000) & 0.000 (0.000) & 0.00 (0.00) & 0.0 (0.0) \\
combined & dense & 0\% & 1.000 (0.000) & 0.000 (0.000) & 0.000 (0.000) & 0.00 (0.00) & 0.0 (0.0) \\
intent links & dense & 40\% & 0.991 (0.001) & 0.009 (0.001) & 0.168 (0.011) & 0.40 (0.01) & 9.4 (0.6) \\
approval links & dense & 40\% & 0.953 (0.002) & 0.047 (0.002) & 0.401 (0.015) & 0.40 (0.01) & 21.1 (0.8) \\
correlation ids & dense & 40\% & 0.840 (0.010) & 0.160 (0.010) & 0.653 (0.013) & 0.97 (0.03) & 49.1 (1.3) \\
timestamps & dense & 40\% & 1.000 (0.000) & 0.000 (0.000) & 0.000 (0.000) & 0.00 (0.00) & 0.0 (0.0) \\
ids plus timestamps & dense & 40\% & 0.648 (0.011) & 0.180 (0.007) & 0.651 (0.012) & 0.86 (0.02) & 47.9 (1.3) \\
source artifacts & sparse & 40\% & 0.332 (0.012) & 0.266 (0.020) & 0.954 (0.010) & 1.41 (0.04) & 145.9 (4.0) \\
source artifacts & dense & 40\% & 0.323 (0.009) & 0.612 (0.009) & 0.953 (0.008) & 2.09 (0.07) & 147.5 (1.8) \\
combined & sparse & 5\% & 0.866 (0.015) & 0.058 (0.012) & 0.318 (0.025) & 0.42 (0.03) & 24.0 (1.9) \\
combined & dense & 5\% & 0.839 (0.008) & 0.145 (0.008) & 0.355 (0.015) & 0.44 (0.02) & 27.0 (1.3) \\
combined & sparse & 20\% & 0.530 (0.017) & 0.137 (0.013) & 0.837 (0.016) & 1.25 (0.04) & 97.5 (3.5) \\
combined & dense & 20\% & 0.502 (0.011) & 0.384 (0.010) & 0.872 (0.009) & 1.57 (0.05) & 105.6 (2.2) \\
combined & sparse & 40\% & 0.231 (0.013) & 0.145 (0.019) & 0.990 (0.006) & 1.43 (0.04) & 186.8 (3.2) \\
combined & dense & 40\% & 0.234 (0.008) & 0.479 (0.019) & 0.992 (0.002) & 1.98 (0.06) & 191.2 (2.0) \\
\bottomrule
\end{tabular}

\end{table*}

\begin{figure*}[t]
\centering
% Simulation results, 2x2 group (two-column width), data from data/series/
\begin{tikzpicture}
\begin{groupplot}[
  group style={group size=2 by 2, horizontal sep=1.55cm, vertical sep=1.45cm},
  edplot,
]
% (a) coverage, dense
\nextgroupplot[title={(a) link coverage (dense)}, ylabel={coverage},
  ymin=0, ymax=1.05, legend pos=south west]
\addplot[serA, mark=*, mark size=1.5pt, thick]
  table[x=p, y=cov, col sep=comma] {data/series/dense_approval_links.csv};
\addlegendentry{approval links}
\addplot[serC, mark=square*, mark size=1.4pt, thick, densely dashed]
  table[x=p, y=cov, col sep=comma] {data/series/dense_correlation_ids.csv};
\addlegendentry{correlation ids}
\addplot[black, mark=pentagon*, mark size=1.6pt, thick]
  table[x=p, y=cov, col sep=comma] {data/series/dense_ids_plus_timestamps.csv};
\addlegendentry{ids $+$ timestamps (pairwise)}
\addplot[serB, mark=triangle*, mark size=1.8pt, thick, dash dot]
  table[x=p, y=cov, col sep=comma] {data/series/dense_source_artifacts.csv};
\addlegendentry{record deletion}
\addplot[serD, mark=diamond*, mark size=1.8pt, thick, densely dotted]
  table[x=p, y=cov, col sep=comma] {data/series/dense_combined.csv};
\addlegendentry{all classes}

% (b) false joins, combined, both profiles
\nextgroupplot[title={(b) false-join rate (all-classes arm)},
  ylabel={false joins / accepted links}, ymin=0, ymax=0.55,
  legend pos=north west]
\addplot[serD, mark=diamond*, mark size=1.8pt, thick]
  table[x=p, y=false, col sep=comma] {data/series/dense_combined.csv};
\addlegendentry{dense}
\addplot[serA, mark=o, mark size=1.5pt, thick, densely dashed]
  table[x=p, y=false, col sep=comma] {data/series/sparse_combined.csv};
\addlegendentry{sparse}

% (c) irrecoverability constructs
\nextgroupplot[title={(c) irrecoverability constructs (dense)},
  ylabel={fraction of chains}, ymin=0, ymax=1.05, legend pos=south east,
  domain=0:40, samples=100]
\addplot[serB, mark=triangle*, mark size=1.8pt, thick, only marks]
  table[x=p, y=recloss, col sep=comma] {data/series/dense_combined.csv};
\addlegendentry{record loss (all classes)}
\addplot[black!60, thick] {1-(1-x/100)^6};
\addlegendentry{$1-(1-p)^{6}$ (consistency check)}
\addplot[serC, mark=square*, mark size=1.4pt, thick, densely dashed]
  table[x=p, y=irrec, col sep=comma] {data/series/dense_correlation_ids.csv};
\addlegendentry{identifiability (key stripping)}

% (d) computed policy-specific ED
\nextgroupplot[title={(d) computed $\ED_\rho$ (dense, declared workload)},
  ylabel={units / configuration}, ymin=0, ymax=240, legend pos=north west]
\addplot[serD, mark=diamond*, mark size=1.8pt, thick]
  table[x=p, y=ed, col sep=comma] {data/series/dense_combined.csv};
\addlegendentry{all classes}
\addplot[serB, mark=triangle*, mark size=1.8pt, thick, dash dot]
  table[x=p, y=ed, col sep=comma] {data/series/dense_source_artifacts.csv};
\addlegendentry{record deletion}
\addplot[serC, mark=square*, mark size=1.4pt, thick, densely dashed]
  table[x=p, y=ed, col sep=comma] {data/series/dense_correlation_ids.csv};
\addlegendentry{correlation ids}
\end{groupplot}
\end{tikzpicture}

\caption{Reconstruction outcomes versus evidence degradation (means over 20
seeds; SDs below marker size except where bars shown). (a)~Link coverage,
dense: key-stripping classes stay near ceiling; the pairwise arm
(ids$+$timestamps) departs from the correlation-ids arm superadditively
(\cref{fnd:superadd}); record deletion and the all-classes arm collapse.
(b)~False-join rate in the all-classes arm: density roughly triples silent
error at fixed degradation. (c)~The two irrecoverability constructs, dense:
record loss in the all-classes arm tracks the source-survival combinatorics
$1-(1-p)^{6}$ (consistency check of the record-loss component, not a
validation of \cref{def:irrec}); identifiability-based irrecoverability under
pure key stripping is large although record loss is identically zero---the
constructs measure different things. (d)~Computed $\ED$ under the declared
workload of \cref{sec:sim-workload}, dense.}
\label{fig:results}
\end{figure*}

\Cref{tab:results} and \cref{fig:results} summarize; released data contain all
cells. Throughout, the largest per-cell coverage SD across the main factorial
is 1.9 points (sparse, all-classes, $p{=}10\%$).

\begin{finding}[Demonstration with removal test: redundancy absorbs key loss;
deletion defeats it]
\label{fnd:redundancy}
At 40\% degradation of any single key class, link coverage remains at
95.3--99.1\% (dense; 98.8--99.9\% sparse) except correlation identifiers
(84.0\% dense, 96.0\% sparse); whole-record deletion at the same rate collapses
coverage to 32--33\% with 95\% of chains suffering record loss. The plateau is
partly a product of the generator's (realistic) redundancy, and the removal arm
shows how much: with redundancy disabled, the correlation-ids arm at
$p{=}40\%$ (dense) falls from 84.0\% to 72.6\% coverage, identifiability-based
irrecoverability rises from 65.3\% to 84.7\% of chains, and computed $\ED$
rises 1.6$\times$ (49.1 to 80.9 units). Redundancy is thus not an artifact to
apologize for but a quantified protective factor---and its protection is
bounded: no redundancy setting rescues deleted records
(\cref{prop:plateau,def:irrec}).
\end{finding}

\begin{finding}[Falsifiable test passed: the Proposition-2 pairwise arm is
superadditive; the all-classes arm is deletion-dominated and subadditive]
\label{fnd:superadd}
The pairwise arm (correlation identifiers $+$ timestamps) instantiates
\cref{prop:superadd}'s premise: neither class alone destroys records (record
loss is 0.000 in every such cell). Timestamps alone have \emph{zero} effect at
every level in both profiles; correlation-ids alone cost 16.0 points of
coverage at $p{=}40\%$ (dense). Their combination costs \textbf{35.2} points
(dense)---more than double the sum---and \textbf{32.6} versus a sum of 4.0
points (sparse), an eight-fold superadditive excess: stripped keys force
heuristic recovery exactly where nulled timestamps have destroyed the
discriminators heuristics need. This is the paper's central interaction
prediction, it could have failed (redundant digest paths might have absorbed
the combination), and it did not. By contrast, the all-classes arm is
\emph{sub}additive in coverage relative to the sum of its five singles (e.g.,
16.1 versus 16.9 points at $p{=}5\%$, dense) and its losses are dominated by
the deletion component (15.0 of those points): once a record is deleted, the
additional loss of its keys or timestamps cannot be charged twice. Interaction
structure, not headline collapse, is where the model earns or loses credit,
and it behaves as \cref{prop:superadd} says: superadditivity arises between
classes that disable \emph{disjoint recovery paths}, and saturates where one
class destroys the nodes outright.
\end{finding}

\begin{finding}[Effort precedes failure; correctness outruns confidence]
\label{fnd:effort}
Under the invariant accounting, baseline effort is zero (every link resolves
exactly). At $p{=}40\%$ correlation-id stripping (dense), 0.97 of 6 links per
chain require heuristic resolution at 146 candidate examinations per chain,
while coverage is still 84\%; approval-link stripping puts 0.40 links per
chain into fallback at 95.3\% coverage. The estate's reconstruction economics
degrade well before its dashboards would: success-rate metrics see 95\%,
while a growing fraction of history is being recovered by search rather than
by lookup. The same cells expose a subtler and operationally important gap:
identifiability-based irrecoverability at that cell is 65.3\% of chains
although the reconstructor's answers are 84\% correct---nearest-in-time
guessing is right far more often than it can \emph{certify}. Reconstructed
history under deficit is thus not merely lossy; it is unauditable even when
correct, which is precisely the confidence-accounting problem
\cref{sec:metrics} flags for ERC reporting.
\end{finding}

\begin{finding}[Falsifiable test passed: density scales silent error, isolated
on a single knob]
\label{fnd:false}
Accepted-but-wrong links reach 14.5\% (sparse) and 47.9\% (dense) of accepted
links in the all-classes arm at $p{=}40\%$, and 61.2\% (dense) under pure
record deletion. Because the sparse/dense profiles bundle three parameters,
the single-knob sweep varies chains-per-service alone (50/100/150; actors and
window fixed) at $p{=}10\%$ all-classes: the false-join rate rises
10.1\% $\to$ 16.1\% $\to$ 19.5\% while coverage and irrecoverability stay flat
($73.6 \to 72.2\%$; $57.0 \to 58.7\%$). Ambient activity density degrades
specifically the \emph{error} channel, as linkage theory
predicts~\cite{fellegi1969}---and as observed in the wild in bug-link datasets,
where heuristic linkage produced systematically biased
corpora~\cite{bird2009,bachmann2010}. False joins are the most dangerous
repayment currency: they produce a confident, wrong account of history.
\end{finding}

\begin{finding}[The two irrecoverability constructs diverge as the model
requires]
\label{fnd:constructs}
Record loss in deletion-bearing arms tracks the source-survival combinatorics
$1-(1-p)^{6}$ within 1.2 points at every level---a consistency check that the
record-loss component measures exactly what it is defined to measure (the
agreement is arithmetically expected, and we claim nothing more from it).
The construct \cref{def:irrec} actually requires is broader, and behaves
differently: under pure key stripping, record loss is identically zero while
identifiability-based irrecoverability reaches 65.3\% (dense) and 29.7\%
(sparse) of chains at $p{=}40\%$---history whose every record survives and
whose connectivity nevertheless cannot be re-established at confidence.
Timestamp nulling alone leaves both constructs at zero (keys never consult
timestamps), the null result that makes the pairwise superadditivity of
\cref{fnd:superadd} interpretable: timestamp evidence is a \emph{conditional}
asset whose value is realized only when keys fail.
\end{finding}

\begin{finding}[$\ED$ computed end to end; its composition is
declaration-sensitive]
\label{fnd:ed}
Under the declared workload, $\ED$ per configuration rises from 0 at baseline
to 191 units (dense, all-classes, $p{=}40\%$), with the full census-to-debt
pipeline of \cref{prop:reduction} exercised in code. Decomposition is
instructive: at every non-trivial level the write-off term dominates the
effort term by an order of magnitude or more (e.g., 51.5 of 53.5 units at
$p{=}10\%$, dense), i.e., under this declaration the economic mass of evidence
debt sits in irrecoverable history, not in reconstruction labor. That
conclusion inverts if the declared penalty $\pi$ is small relative to effort
costs---which is not a weakness of the construct but its content:
\cref{def:debt} makes the workload-and-penalty judgment an explicit input, and
the simulation shows exactly how the judgment propagates
(\cref{sec:disc-relativity}).
\end{finding}

\subsection{Sensitivity of the Apparatus Itself}
\label{sec:sim-sensitivity}
Beyond the redundancy-removal and density-sweep arms already reported:
narrowing the heuristic window to 24 hours---\emph{below} the 36-hour maximum
true inter-record gap, so that some true parents are structurally outside the
search---leaves coverage nearly unchanged (83.4\% versus 84.0\% at the
correlation-ids/dense/40\% cell) while \emph{reducing} identifiability-based
irrecoverability (57.7\% versus 65.3\%): a narrower window excludes some true
parents but disambiguates many more candidate sets. Window choice thus trades
reach against confidence rather than simply scaling results. Making the
ticket filer differ from the commit author in 30\% of chains (breaking an
actor-match assumption the generator otherwise makes perfectly reliable)
lowers intent-link coverage at $p{=}40\%$ (dense) from 99.1\% to 97.5\% and
raises identifiability-based irrecoverability from 16.8\% to 23.4\%---a
measurable, modest degradation showing which headline numbers lean on that
assumption. We report these arms so that readers can see the apparatus move,
and the artifact accepts further policies by substitution.

\subsection{Threats to Validity}
\label{sec:sim-threats}
\emph{Construct}: fallback-link counts and candidate examinations are
implementation-invariant but still algorithmic; human reconstruction adds
search, judgment, and interruption costs---the field drills of \cref{sec:field}
calibrate the mapping, and the declared unit costs of \cref{sec:sim-workload}
are placeholders whose only role is demonstrating the aggregation.
\emph{Internal}: uniform-random degradation is a simplification; real decay is
correlated (a migration deletes a store wholesale)---correlated-decay arms are
a supported extension. The generator and reconstructor share authorship with
the model; the removal and sensitivity arms are our defense, and the labeled
separation of demonstrations from falsifiable tests is the honest reading of
what each result can carry. \emph{External}: the estate is synthetic;
magnitudes do not transfer, and only the interaction structure---plateaus with
quantified redundancy dependence, pairwise superadditivity, effort-before-%
failure, density-scaled error---is claimed, each tied to a model mechanism.
Real-world evidence that the phenomenon and its error channel exist at scale
is, however, not absent from the literature: the mining-software-repositories
community measured missing bug-fix links and the bias of heuristically
reconstructed histories over a decade ago~\cite{bachmann2010,bird2009}.
\emph{Conclusion}: per-cell SDs are $\le$1.9 coverage points; the differences
supporting \cref{fnd:redundancy,fnd:superadd,fnd:false,fnd:constructs} exceed
that by an order of magnitude and are stable across seeds.
